\begin{abstract}

在数字化教育快速发展的背景下，课程电子化与教学互联网化日益普及，本科程序设计课程中的一体化在线评测平台便是数字化浪潮中的重要产物。然而，现有平台在实际教学中存在着一常见问题：学生提交代码后只能看到最终判题结论，例如“答案错误”或“编译错误”，无法判断失败是发生在排队、编译还是运行阶段；同时平台开发方又不得不在安全隔离、并发处理能力与反馈可解释性之间做出工程妥协。针对上述问题，本文设计并实现了一个面向本科程序设计课程的一体化在线评测平台。该系统将客观判题作为主干，AI分析只做辅助，两者之间通过数据交换联系，但不共用流程，从而保证判题结果的稳定性不受AI模块波动的影响。

本系统的判题流程将提交受理和执行分开处理。用户提交代码时，系统先创建一条状态为排队中的记录，同时投递一个 Inngest 事件。后台函数准备 Docker 容器、上传源码、完成编译、按测试用例逐一运行并清理现场。执行过程放在受控的 Docker 沙箱内，并设置了内存上限和网络隔离。状态变化通过 Soketi 实时推送到前端，避免轮询。为便于排查问题，数据层没有只存最终结果，而是拆成三个模型：Submission 存放提交的基本信息，Judge 记录本次判题的整体状态，JudgeRun 保存每个测试用例的输入、输出、消耗时间和内存。

本系统还围绕代码提交的前后环节给出教学方面的支持，提交之前，系统利用 Monaco 编辑器还有容器化 Clangd 服务给出代码补全、语法诊断还有函数签名提示功能；提交之后，AI 模块会异步生成复杂度、代码风格、可读性、运行效率以及正确性等内容的分析结果，还会把这些分析结果跟客观判题结果分开储存。验证结果表明，系统可以稳定处理正确通过、编译错误、答案错误、运行时错误、超时以及内存超限这些常见情况，还能借助资源限制和网络隔离阻挡恶意提交。跟只给出最终判题结论的传统平台相比，本系统更符合本科程序设计教学的实际需求。

\end{abstract}

\begin{abstractEn}

This paper implements an integrated online evaluation platform for undergraduate programming teaching scenarios. The starting point of the topic selection comes from the actual teaching pain points : students often can only see the final judgment state, and it is difficult to judge whether the failure occurs in the queuing, compiling or running stage, while the platform side needs to make engineering choices between security isolation, concurrent processing and feedback interpretability. In response to these problems, this paper takes objective judgment as the core process, and AI ability as an independent teaching assistant process, which is decoupled from the judgment process to ensure the stability and certainty of the judgment process.

This paper adopts the decoupling scheme of submission acceptance and asynchronous execution in the judgment process. The submission entry first creates a record with a state of QD and delivers the Inngest event. The background function completes the container preparation, source code upload, compilation, run and clean up one by one according to the test case. The  execution phase runs in a controlled Docker sandbox, and combines memory constraints and network isolation to guarantee execution boundaries. The state change is transmitted back to the front end in real time through the Soketi channel to reduce the polling overhead. In order to improve the traceability, the data layer divides the judgment information into three layers : Submission, Judge and JudgeRun, which respectively carry the submission summary, the main state of the judgment and the execution track of each test case, so as to completely record the whole process of the judgment process from acceptance to final state.

On the teaching support side, the platform covers support capabilities in two stages, before and after submission. Before submission, the front-end workspace integrates Monaco with Language Server Protocol to provide completion, diagnostics, and signature help through containerized Clangd. After submission, the AI module asynchronously generates multi-dimensional analysis results, including complexity, comprehensive score, code style, readability, efficiency, and correctness, and provides executable improvement suggestions. The analysis results are stored separately from objective judgment results to avoid the interpretation module affecting the certainty of the ruling. The system supports Chinese and English interface and topic switching, and the AI analysis text can output Chinese or English feedback according to the current language environment. Verification results show that the platform can stably complete AC, CE, WA, RE, TLE, and MLE core state determination, and the asynchronous link and real-time state push can work continuously. In response to malicious samples such as \texttt{fork()} abuse, oversized memory requests, and network outreach requests, the system can also block risky behavior under resource restriction and network isolation strategies. Compared with the traditional teaching OJ represented by UOJ, which mainly focuses on the final verdict, this system provides integrated capabilities of semantic diagnosis before submission, real-time state visualization during submission, and multi-dimensional interpretation feedback after submission while maintaining the stability of objective judgment, reflecting system features and engineering innovation oriented to teaching process support.

\end{abstractEn}
